\centerline{\bf \Huge Четность}


\begin{flushright}
        \scriptsize{— Патрик, хорошая идея!\\
                    — Это моё призвание.\\
                    — Что? Хорошие идеи?\\
                    — Нет, называться Патриком.}
\end{flushright}

\problem{0.1} Покажите, что любое целое число, оканчивающиеся на 0, 2, 4, 6 и 8, делится на 2.

\problem{0.2} 
Покажите, что сумма или разность двух целых чисел одинаковой четности четна

\problem{0.3} 
Покажите, что сумма четного числа нечетных слагаемых четна

\problem{0.4} Сумма трех чисел нечетна. Сколько нечетных слагаемых может быть?

\problem{1} Можно ли разменять купюру в 25 рублей на десять купюр по 1, 3 и 5 рублей?

\problem{2} 
Сайгак прыгает по прямой: первый раз на 1 см, второй - на 2 см и т. д. Может ли он через 25 прыжков вернуться на старое место?


\problem{3}
Не вычисляя суммы $1 + 2 + 3 + \dots + 2024 $ определите ее четность.

\problem{4} На столе стоят шесть столбиков монет. В первом столбике одна монета, во втором - две, в третьем - три, ..., в шестом - шесть. Разрешается на любые два столбика положить по монете. а) Можно ли за несколько таких операций сделать все столбики одинаковыми? б) А что если в последнем столбике будет 7 монет?

\problem{5} В стоэтажной кибитке испортился лифт. Теперь в нем работают только две кнопки. При нажатии на первую лифт поднимется на 8 этажей, при нажатии на вторую - опускается на 6 этажей. Можно ли попасть с первого этажа а) на 98-й, б) на 95-й.  

\problem{6}
Может ли конь пройти с поля а1 на поле h8 (пра-вое верхнее), побывав по дороге на всех полях ровно по одному разу?
90


